\documentclass[11pt,reqno]{amsart}
\usepackage[OT1]{fontenc}
\usepackage{amsmath,amssymb,mathtools}
\usepackage{microtype}
\usepackage{needspace}
\usepackage[hidelinks,unicode]{hyperref}
\hypersetup{pdftitle={A rational counterexample to the reciprocal-logarithmic assertion in Erdos problem 143},pdfauthor={Mattia Apicella},pdfsubject={A formally verified construction; external expert review pending}}
\newtheorem{theorem}{Theorem}[section]
\newtheorem{proposition}[theorem]{Proposition}
\newtheorem{lemma}[theorem]{Lemma}
\newtheorem{corollary}[theorem]{Corollary}
\theoremstyle{definition}
\newtheorem{definition}[theorem]{Definition}
\theoremstyle{remark}
\newtheorem{remark}[theorem]{Remark}
\numberwithin{equation}{section}
\newcommand{\N}{\mathbb N}
\newcommand{\Z}{\mathbb Z}
\newcommand{\Q}{\mathbb Q}
\newcommand{\R}{\mathbb R}
\newcommand{\C}{\mathbb C}
\newcommand{\Om}{\Omega}
\newcommand{\wt}{\mathrm w}
\newcommand{\dens}{\mathrm d}
\newcommand{\lcm}{\operatorname{lcm}}
\newcommand{\cR}{\mathcal R}
\newcommand{\cA}{\mathcal A}
\newcommand{\abs}[1]{\lvert #1\rvert}
\newcommand{\set}[1]{\{#1\}}
\title[A rational counterexample for Erd\H{o}s problem 143]{A rational counterexample to the reciprocal-logarithmic assertion in Erd\H{o}s problem 143}
\author{Mattia Apicella}
\date{October 2, 2026; revised October 23, 2026}
\subjclass[2020]{Primary 11J25; Secondary 11B05, 11N37, 68V20}
\keywords{Integer dilation approximation, primitive sets, periodic sets, Dirichlet characters}
\begin{document}
\begin{abstract}
We construct a countably infinite set $A\subset\mathbb Q\cap[2,\infty)$
such that $|ka-b|\geq1$ for all distinct $a,b\in A$ and all positive
integers $k$, while $\sum_{a\in A}1/(a\log a)=\infty$. This gives a
negative answer to the reciprocal-logarithmic convergence question in
Erd\H{o}s problem 143, distinct from its logarithmic-density question.
The construction combines a primitive-set selection lemma for periodic
sets with rational blocks equipped with private primes. At the beginning
of each round, a new integer grid makes the preceding configuration
exclude only a small proportion of labels. Finitely many new blocks then
add a uniformly positive amount of weight without moving earlier points. The selection
lemma follows from Euler-product estimates, Dirichlet characters and
Abel averaging at a fixed modulus. Every constructed point lies within
$m^{-12}/2$ to the right of an integer $m$. A companion Lean~4
formalization checks the analytic lemma, the finite extensions and the
complete existence theorem.
\end{abstract}
\maketitle

\section{Introduction}
\label{sec:introduction}

Write $\N=\{1,2,\ldots\}$ and let $\log$ denote the natural logarithm.
For $x>1$ and for a finite or countable set $B\subset(1,\infty)$, put
\[
 \wt(x)=\frac1{x\log x},\qquad
 W(B)=\sum_{x\in B}\wt(x)\in[0,\infty].
\]
The sum is understood as the supremum of the sums over finite subsets.

\begin{definition}
A set $A\subset(1,\infty)$ is \emph{dilation-separated} if
\begin{equation}\label{eq:original}
 \abs{ka-b}\geq1
 \quad\text{for all }a,b\in A,\ a\ne b,\quad
 \text{and all }k\in\N.
\end{equation}
The quantification is over ordered pairs, so both directions are required.
A set $P\subset\N\setminus\{1\}$ is \emph{primitive} if no element of $P$
divides a distinct element of $P$.
\end{definition}

Taking $k=1$ in~\eqref{eq:original} shows that distinct points are at
least one apart. Hence every bounded interval contains only finitely
many points of a dilation-separated set, and every such set is at most
countable.

For a set of integers greater than one, condition~\eqref{eq:original}
is equivalent to primitivity: an integer distance less than one must be
zero, so the excluded event is precisely divisibility between distinct
elements. Erd\H{o}s proved that every primitive set of integers has finite
reciprocal-logarithmic weight~\cite{Erdos1935}. The analogous question
for real sets is subtler, since exact divisibility is replaced by
uniform separation from integer dilates.

In problem II on the final page of~\cite{Erdos1948}, Erd\H{o}s asks whether
condition~\eqref{eq:original} forces both convergence of $W(A)$ and
vanishing logarithmic density. These are the two questions in the problem
now numbered 143 in~\cite{Bloom143}. Koukoulopoulos, Lamzouri, and
Lichtman~\cite[Theorem 1]{KLL2025} established the latter assertion:
every dilation-separated set satisfies
\begin{equation}\label{eq:logdensity}
 \sum_{a\in A\cap[1,x]}\frac1a=o(\log x).
\end{equation}
Their theorem gives approximate integer-dilation relations under positive
upper logarithmic density. The theorem below shows that divergence of
$W(A)$ alone does not force such relations, even for rational sets whose
points lie extremely close to integers. In particular, the two proposed
sparsity conclusions have different answers; the counterexample is
consistent with~\eqref{eq:logdensity}.

\begin{theorem}\label{thm:main}
There exists a countably infinite set $A\subset\Q\cap[2,\infty)$ with
no finite accumulation point such that \eqref{eq:original} holds and
\begin{equation}\label{eq:divergence}
 W(A)=\infty.
\end{equation}
Every $a\in A$ can be written as
\begin{equation}\label{eq:near-integer}
 a=m+\eta,
 \qquad m\in\N,\quad 0<\eta\leq\frac12m^{-12}.
\end{equation}
More precisely, the construction produces one nested sequence of finite
dilation-separated sets $F_1\subset F_2\subset\cdots$ whose union is $A$ and for which
\begin{equation}\label{eq:prefix-weight}
 W(F_R)\geq\frac{R}{576}\qquad(R\in\N).
\end{equation}
\end{theorem}

The main difficulty is to accumulate infinite weight while maintaining
separation from every integer multiple of every earlier point. A small
perturbation of integer primitive sets alone does not give this global
compatibility. We use finite rational blocks with a private-prime
invariant. At the beginning of each round this invariant supplies a new
grid on which all earlier points exclude only a small proportion of
labels. The grid changes, but the earlier points do not. Periodic
primitive-set selection then supplies enough new labels to obtain a
uniform gain in weight. The constants are sufficient for this construction;
no optimality is asserted for the gain or the exponent in
\eqref{eq:near-integer}.

The following sections give the argument in full. The companion formal
proof has no additional analytic or construction hypothesis; its scope
and reproduction instructions appear in Section~\ref{sec:verification}.

\subsection{Outline and notation}
For $n\in\N$, let $\Om(n)$ be the number of prime factors of $n$, counted
with multiplicity; in particular, $\Om(1)=0$. Every subset of
\[
 \mathcal L_r=\{n\in\N:\Om(n)=r\},\qquad r\geq1,
\]
is primitive. Indeed, $n\mid n'$ with $n\ne n'$ implies
$\Om(n')=\Om(n)+\Om(n'/n)>\Om(n)$.

A set $E\subset\N$ is \emph{periodic of period $L$} if membership depends
only on the residue modulo $L$, where $L\in\N$. Its density is
\[
 \dens(E)=\frac{\#\{b\in\Z/L\Z:b\text{ is admitted by }E\}}{L}.
\]
This definition is independent of the chosen period. All density
calculations below take place on a finite common period. In particular,
the union bound for finitely many periodic sets is simply the union bound
for finite subsets of a residue ring.

The analytic part supplies a finite primitive subset of weight greater
than $1/6$ in every periodic set of density greater than $2/3$, above any
fixed threshold. The geometric part uses it to construct a block of
weight greater than $1/(24D)$ while spending at most $3/D$ of the density
of labels on a fixed grid $D\N$. A round contains $\lfloor D/12\rfloor$
blocks. Private primes make it possible to restart the next round on a
new grid while preserving the entire preceding configuration.

\subsection*{AI Declaration}
OpenAI Codex was used during the research and preparation of this article.
GPT-5.6 Sol was used in the initial exploration of a positive answer to the
convergence question, and Astra was used in developing the counterexample.
AI assistance contributed to mathematical
exploration, proof development, Lean formalization, and exposition,
including the editorial revision of this manuscript. The historical model
identifiers are reported as retained by the author; a more specific
version identifier for Astra is not recorded here. The author takes
responsibility for the content of the article.

\subsection*{Mathematical software and supporting material}
The companion project uses Lean~4.24.0 and a pinned revision of mathlib.
The predicate named \path{Original} in the code is precisely the
dilation-separation condition~\eqref{eq:original}.
Its source files, exact theorem statements and reproduction instructions
are supplied as supplementary material; the repository is
\url{https://github.com/mapice/erdos-143-candidate-disproof}.
The scope of the formal checks is described in
Section~\ref{sec:verification}.

\section{Primitive weight in a periodic set}
\label{sec:analytic}

The classical analytic inputs in this section are the absolutely
convergent Euler products for $\zeta(s)$ and Dirichlet $L$-functions for
$\Re s>1$, character orthogonality for a fixed modulus, and continuity of
$L(s,\chi)$ at $s=1$ for nonprincipal $\chi$; see \cite[Sections 2 and 6]{Tao2014}.
Non-vanishing at $s=1$ is not needed. We use an Abel average to find
arbitrarily high levels with sufficient weight, without requiring a
lower bound for every sufficiently high individual level.

\begin{theorem}[Periodic selection]\label{thm:selection}
Let $E\subset\N$ be periodic, with $\dens(E)>2/3$. For every integer
$N\geq2$, there is a finite primitive set
$P\subset E\cap[N,\infty)$ such that $W(P)>1/6$.
\end{theorem}

For a fixed modulus $L\geq1$ and residue $b\pmod L$, define
\begin{equation}\label{eq:abel-residue-mass}
 A_{L,b}(u)=\sum_{\substack{n\geq2\\n\equiv b\pmod L}}
                 u^{\Om(n)}\wt(n),\qquad 0\leq u<1.
\end{equation}
We will prove convergence of this series and the sufficient estimate
\begin{equation}\label{eq:residue-abel}
 \liminf_{u\uparrow1}(1-u)A_{L,b}(u)\geq\frac1{4L}.
\end{equation}
The modulus and residue are fixed before the limit is taken; no
uniformity in $L$ is required. This estimate will supply arbitrarily
high primitive levels whose weight exceeds $1/6$ after summation over
the admitted residue classes.

\subsection{Uniform Euler-factor estimates}
For a Dirichlet character $\chi$ modulo $L$, $0\leq u\leq1$, and $t>0$,
put
\begin{equation}\label{eq:generating-product}
 G_\chi(u,t)=\prod_p(1-u\chi(p)p^{-1-t})^{-1}
          =\sum_{n\geq1}\frac{u^{\Om(n)}\chi(n)}{n^{1+t}}.
\end{equation}
Products indexed by $p$ are over primes. Absolute convergence follows
from $|\chi(n)|\leq1$ and $u^{\Om(n)}\leq1$, or directly from finite
Euler products and the summability of $\sum_p p^{-1-t}$. The $n=1$
term is $1$, including when $u=0$.

\begin{lemma}\label{lem:euler-comparisons}
For $0\leq u\leq1$ and $t>0$,
\begin{equation}\label{eq:twisted-product-bound}
 |G_\chi(u,t)|^2\leq2|L(1+t,\chi)|^{2u}.
\end{equation}
For the principal character $\chi_0$ modulo $L$, the product is real
and positive and satisfies
\begin{equation}\label{eq:principal-product-bound}
 G_{\chi_0}(u,t)\geq\bigl(L(1+t,\chi_0)/2\bigr)^u.
\end{equation}
\end{lemma}

\begin{proof}
For $z\in\C$ with $|z|\leq1/2$, set
$A=|1-z|^2$ and $S=(1-u)+uA$. The identity
\[
 |1-uz|^2=S-u(1-u)|z|^2
\]
and the bounds $S\geq1/4$ and $u(1-u)\leq1/4$ give
\[
 |1-uz|^2\geq S(1-|z|^2)\geq A^u(1-|z|^2),
\]
where the last step is weighted arithmetic--geometric mean. Inverting,
\begin{equation}\label{eq:single-factor}
 |(1-uz)^{-1}|^2
 \leq |(1-z)^{-1}|^{2u}(1-|z|^2)^{-1}.
\end{equation}
The quadratic corrections have a uniform bound. For any finite set
$T$ of primes and $M\geq\max(T\cup\{2\})$,
\begin{equation}\label{eq:correction-product}
 \prod_{p\in T}(1-p^{-2})^{-1}
 \leq\prod_{n=2}^M(1-n^{-2})^{-1}
 =\frac{2M}{M+1}<2.
\end{equation}
Apply \eqref{eq:single-factor} with $z=\chi(p)p^{-1-t}$ and pass to
the convergent Euler products to obtain
\eqref{eq:twisted-product-bound}.

For $0\leq x<1$, concavity and an elementary rearrangement give
\[
 (1+x)^u\leq1+ux\leq(1-ux)^{-1}.
\]
For the principal character use $x=p^{-1-t}$ when $p\nmid L$ and
$x=0$ otherwise. Since $1+x=(1-x)^{-1}(1-x^2)$,
\eqref{eq:correction-product} implies
\[
 G_{\chi_0}(u,t)\geq
 \left(\prod_{p\nmid L}(1+p^{-1-t})\right)^u
 \geq\bigl(L(1+t,\chi_0)/2\bigr)^u.
\]
The same argument applies to finite products first; convergence then
justifies the displayed passage to infinite products.
\end{proof}

Write $G_1$ for the unrestricted product, corresponding to modulus one,
and put
\[
 \kappa_L=\prod_{p\mid L}(1-p^{-1})=\frac{\varphi(L)}L,
 \qquad \varphi(1)=1.
\]

\begin{lemma}\label{lem:generating-bounds}
For a fixed nonprincipal character $\chi$,
\begin{equation}\label{eq:nonprincipal-local}
 |G_\chi(u,t)|\leq C_{L,\chi}
 \qquad(0\leq u\leq1,\ 0<t\leq1).
\end{equation}
There is $\varepsilon\in(0,1]$, depending on $L$, such that
\begin{equation}\label{eq:principal-local}
 G_{\chi_0}(u,t)\geq\left(\frac{\kappa_L}{4t}\right)^u
 \qquad(0\leq u\leq1,\ 0<t\leq\varepsilon).
\end{equation}
The unrestricted product satisfies
\begin{equation}\label{eq:unrestricted-local}
 G_1(u,t)\leq4t^{-u}
 \qquad(0\leq u\leq1,\ 0<t\leq1).
\end{equation}
Uniformly in all characters and $0\leq u\leq1$,
\begin{equation}\label{eq:generating-tail}
 |G_\chi(u,t)-1|\leq C2^{-t}\qquad(t\geq1),
\end{equation}
with an absolute constant $C$.
\end{lemma}

\begin{proof}
For nonprincipal $\chi$, continuity on the compact real segment $[1,2]$
gives a bound $B\geq1$ with $|L(s,\chi)|\leq B$ there. For $0<u\leq1$,
equation~\eqref{eq:twisted-product-bound} gives
$|G_\chi(u,t)|^2\leq2B^{2u}\leq2B^2$, while $G_\chi(0,t)=1$.
This proves \eqref{eq:nonprincipal-local} with $C_{L,\chi}=2B$.
The principal Euler product and the elementary estimates
\[
 \frac1t\leq\zeta(1+t)\leq1+\frac1t
\]
give $tL(1+t,\chi_0)\to\kappa_L$ as $t\downarrow0$.
In particular, $L(1+t,\chi_0)\geq\kappa_L/(2t)$ for small $t$.
Equation~\eqref{eq:principal-product-bound} proves
\eqref{eq:principal-local}. For modulus one,
\eqref{eq:twisted-product-bound} gives, when $0<t\leq1$,
\[
 G_1(u,t)\leq\sqrt2\,\zeta(1+t)^u
 \leq\sqrt2(2/t)^u\leq4t^{-u}.
\]
Finally, by the absolutely convergent series in
\eqref{eq:generating-product},
\[
 |G_\chi(u,t)-1|\leq\sum_{n\geq2}n^{-1-t}
 \leq2^{-(t-1)}\sum_{n\geq2}n^{-2}\qquad(t\geq1).
\]
This proves \eqref{eq:generating-tail}.
\end{proof}

\subsection{Integration and the character averages}

\begin{lemma}\label{lem:integrability}
For $0\leq u<1$, the series $\sum_{n\geq2}u^{\Om(n)}\wt(n)$
converges. For every character $\chi$,
\begin{equation}\label{eq:weighted-generating}
 J_\chi(u):=\int_0^\infty(G_\chi(u,t)-1)\,dt
     =\sum_{n\geq2}\chi(n)u^{\Om(n)}\wt(n),
\end{equation}
with an absolutely convergent integral and series. Moreover,
\begin{equation}\label{eq:moment-bound}
 C_2:=\sum_{n\geq2}\frac1{n(\log n)^2}<\infty.
\end{equation}
\end{lemma}

\begin{proof}
The unrestricted nonnegative integrand $G_1(u,t)-1$ is integrable:
near zero use \eqref{eq:unrestricted-local} and $u<1$, and near infinity
use \eqref{eq:generating-tail}. Tonelli's theorem and the elementary
identity
\[
 \int_0^\infty n^{-1-t}\,dt=\frac1{n\log n}\qquad(n\geq2)
\]
therefore prove convergence of the unrestricted weighted series and
\eqref{eq:weighted-generating} for the trivial character. For a general
character, the sum of the integrals of the absolute values is bounded by
this same convergent series. Termwise integration is consequently
justified and proves \eqref{eq:weighted-generating}.

For the last assertion, $G_1(1,t)=\zeta(1+t)$. Thus
$t(\zeta(1+t)-1)$ is bounded near zero by the elementary zeta estimates,
and decays at infinity by \eqref{eq:generating-tail}. A further
application of Tonelli, now using
\[
 \int_0^\infty t n^{-1-t}\,dt=\frac1{n(\log n)^2},
\]
proves \eqref{eq:moment-bound}.
\end{proof}

\begin{lemma}\label{lem:character-abel}
For each nonprincipal $\chi$,
\begin{equation}\label{eq:nonprincipal-abel}
 (1-u)J_\chi(u)\longrightarrow0\qquad(u\uparrow1).
\end{equation}
For the principal character,
\begin{equation}\label{eq:principal-abel}
 \liminf_{u\uparrow1}(1-u)J_{\chi_0}(u)\geq\frac{\kappa_L}{4}.
\end{equation}
\end{lemma}

\begin{proof}
For nonprincipal $\chi$, equations~\eqref{eq:nonprincipal-local} and
\eqref{eq:generating-tail} provide an integrable majorant independent of
$0\leq u<1$ for $|G_\chi(u,t)-1|$. Thus $J_\chi(u)$ stays bounded,
which proves \eqref{eq:nonprincipal-abel}.
For $\chi_0$, every term of $G_{\chi_0}(u,t)-1$ is nonnegative.
Integrating the lower bound \eqref{eq:principal-local} on
$(0,\varepsilon)$ gives
\[
 J_{\chi_0}(u)\geq
 \frac{(\kappa_L/4)^u\varepsilon^{1-u}}{1-u}-\varepsilon.
\]
After multiplication by $1-u$, the right side tends to $\kappa_L/4$.
\end{proof}

\subsection{Every residue class}

\begin{proposition}\label{prop:residues}
Estimate \eqref{eq:residue-abel} holds for every fixed $L\geq1$ and every
residue $b\pmod L$, including noninvertible residues.
\end{proposition}

\begin{proof}
For an invertible residue, character orthogonality and the absolute
convergence in Lemma~\ref{lem:integrability} give
\begin{equation}\label{eq:orthogonality}
 A_{L,b}(u)=\frac1{\varphi(L)}
       \sum_{\chi\bmod L}\overline{\chi(b)}J_\chi(u).
\end{equation}
The principal term has normalized lower limit at least
$\kappa_L/(4\varphi(L))=1/(4L)$. The nonprincipal terms, after
multiplication by $1-u$, tend to zero by
Lemma~\ref{lem:character-abel}. Taking real parts in
\eqref{eq:orthogonality} proves the assertion. For $L=1$, the same
argument has only the principal term.

For a general residue use its representative $b\in\{0,\ldots,L-1\}$
and put $g=\gcd(b,L)$, $M=L/g$, and $c=b/g$, with $\gcd(0,L)=L$.
The residue $c\pmod M$ is invertible, also when $M=1$ with its unique
class. Multiplication $m\mapsto gm$ sends the integers $m\geq2$ in
that class injectively into the integers $n\geq2$ in the original
class, and $\Om(gm)=\Om(g)+\Om(m)$. For $m\geq2$,
\begin{equation}\label{eq:log-shift}
 0\leq\frac1{m\log m}-\frac1{m(\log m+\log g)}
 \leq\frac{\log g}{m(\log m)^2}.
\end{equation}
Using $0\leq u^{\Om(m)}\leq1$, $u^{\Om(g)}\leq1$, and
Lemma~\ref{lem:integrability}, we get
\begin{equation}\label{eq:residue-scaling}
 A_{L,b}(u)\geq\frac{u^{\Om(g)}}g A_{M,c}(u)
                    -\frac{\log g}{g}C_2.
\end{equation}
An omitted target term corresponding to $m=1$, if present, has $n=g\geq2$
and is nonnegative, so does not affect the lower bound. Multiplication
by $1-u$ makes the error term vanish.
The first coefficient tends to $1/g>0$, and the coprime case supplies
the lower bound $1/(4M)>0$. The definition of lower limit now gives
$1/(4gM)=1/(4L)$. No coprimality between $g$ and $m$ is required.
\end{proof}

\Needspace{6\baselineskip}
\subsection{Heavy levels and finite extraction}

\begin{proof}[Proof of Theorem~\ref{thm:selection}]
Fix a period $L$ and its set $S$ of admitted residue classes. Sum
Proposition~\ref{prop:residues} over this finite set to obtain
\begin{equation}\label{eq:periodic-abel}
 \liminf_{u\uparrow1}(1-u)
   \sum_{\substack{n\geq2\\n\in E}}u^{\Om(n)}\wt(n)
 \geq\frac{|S|}{4L}=\frac{\dens(E)}4>\frac16.
\end{equation}
For $r\geq0$, put
\[
 f_r=\sum_{\substack{n\geq2\\n\in E\\\Om(n)=r}}\wt(n).
\]
Each $f_r$ is finite: the convergent weighted series at $u=1/2$ bounds
it after multiplication by $2^r$. Regrouping the nonnegative summable
series gives, for $0\leq u<1$,
\begin{equation}\label{eq:layer-decomposition}
 \sum_{\substack{n\geq2\\n\in E}}u^{\Om(n)}\wt(n)
 =\sum_{r\geq0}u^r f_r.
\end{equation}
If $f_r\leq1/6$ for all $r\geq R$, then
\[
 (1-u)\sum_{r\geq0}u^r f_r
 \leq\frac16+(1-u)\sum_{r<R}f_r.
\]
Its upper limit would be at most $1/6$, contradicting
\eqref{eq:periodic-abel}. Thus there are arbitrarily high $r$ with
$f_r>1/6$. Choose such an $r\geq N\geq2$. Every integer on that level
is at least $2^r\geq r\geq N$, and the level is primitive. Since its
weight is a convergent positive series exceeding $1/6$, a finite
truncation still exceeds $1/6$. This truncation is the required set $P$.
\end{proof}

\section{Finite configurations and private primes}
\label{sec:invariants}

We index blocks by $i=1,2,\ldots$ in their order of insertion. A block
has data
\begin{equation}\label{eq:block-data}
 \begin{gathered}
 D_i\in\N,\quad D_i\geq48,\quad
 P_i\subset\N\setminus\{1\}\text{ finite, nonempty, primitive},\\
 H_i,G_i\text{ primes},\qquad
 p_i=H_iG_i,\quad q_i=p_i-1,\quad c_i=\frac{p_i}{q_i},\\
 A_i=\{a_{i,n}:n\in P_i\},\qquad
 a_{i,n}=c_iD_i n=\frac{H_iN_{i,n}}{q_i},\quad
 N_{i,n}=G_iD_i n.
 \end{gathered}
\end{equation}
The prime $H_i$ is called the \emph{private prime} of the block.
Since $q_i\equiv-1\pmod{H_i}$, the factor $H_i$ in the numerator cannot
cancel against the denominator. The invariant below prevents this prime
from dividing the nonprivate numerator parts of any block.

\begin{definition}\label{def:admissible}
A finite list of blocks is an \emph{admissible configuration} if, writing
$F=\bigcup_{i=1}^t A_i$, the following hold:
\begin{enumerate}
\item[(I1)] $F$ is dilation-separated, and every point of $A_i$ exceeds every point
of $\bigcup_{j<i}A_j$ by more than $1$.
\item[(I2)] The primes $H_i$ are distinct and greater than $2$, and
\begin{equation}\label{eq:private-invariant}
 H_j\nmid N_{i,n}\qquad
 (1\leq i,j\leq t,\ n\in P_i).
\end{equation}
\item[(I3)] Every point $a_{i,n}$ satisfies
$a_{i,n}/H_i>4$ and \eqref{eq:near-integer} with $m=D_i n$.
\item[(I4)] The private primes satisfy the strict finite budget
\begin{equation}\label{eq:reciprocal-primes}
 \sum_{i=1}^t\frac1{H_i}<\frac1{100}.
\end{equation}
\end{enumerate}
The empty list is admissible, with $F=\varnothing$. We write
$M(F)=\max F$ when $F\ne\varnothing$, and $M(\varnothing)=0$.
\end{definition}

All invariants concern only data that have already been selected. A new
private prime is chosen large enough to preserve the remaining positive
budget; no prescribed rate of prime growth is needed.


\section{A grid compatible with the preceding configuration}
\label{sec:grid}

\begin{lemma}\label{lem:grid}
Given an admissible configuration with point set $F_0$, there exists an
integer $D\geq48$ such that
\begin{equation}\label{eq:grid-choice}
 \begin{gathered}
 N_{i,n}\mid D\quad\text{for all preceding points},\\
 (D,H_i)=1\quad\text{for all preceding private primes}.
 \end{gathered}
\end{equation}
For this $D$, if $a=a_{i,n}\in F_0$, $\ell\in\N$, and $H_i\nmid\ell$,
then
\begin{equation}\label{eq:grid-separation}
 |ka-D\ell|>4\qquad\text{for every }k\in\Z.
\end{equation}
If instead $H_i\mid\ell$, then $D\ell$ is a positive integer multiple
of $a$.
\end{lemma}

\begin{proof}
For $F_0=\varnothing$, take $D=48$. Otherwise let
$Q$ be the product of the distinct positive integers $N_{i,n}$ stored
in $F_0$. By
\eqref{eq:private-invariant}, $Q$ is coprime to every preceding private
prime. All these primes are odd. Thus $D=2^vQ$, with $v\geq0$ sufficiently
large, satisfies \eqref{eq:grid-choice}.

Fix a preceding point $a=H_iN_{i,n}/q_i$ and put $t_a=D/N_{i,n}$.
This is a positive integer. Since $H_i\nmid D$ and
$q_i\equiv-1\pmod {H_i}$, we have $(t_aq_i,H_i)=1$. The exact identity
\begin{equation}\label{eq:integer-gap}
 |ka-D\ell|=\frac a{H_i}|kH_i-t_aq_i\ell|
\end{equation}
holds for every $k\in\Z$. If $H_i\nmid\ell$, its integer factor is
nonzero, so \eqref{eq:grid-separation} follows from (I3). If
$H_i\mid\ell$, then $D\ell/a=t_aq_i\ell/H_i$ is a positive integer.
The proof uses the displayed representation of $a$ and does not require
that representation to be reduced.
\end{proof}

For a finite rational set $F\subset(0,\infty)$ and an integer $D\geq1$,
define the set of admissible labels
\begin{equation}\label{eq:free-labels}
 \cR(F,D)=\{\ell\in\N:ka-D\ell\notin(-1,1]
             \text{ for every }a\in F\text{ and every }k\in\Z\}.
\end{equation}
The half-open interval is chosen for the later positive perturbation.

\begin{lemma}\label{lem:periodicity}
The set $\cR(F,D)$ is periodic. For $F=F_0$ and $D$ as in
Lemma~\ref{lem:grid},
\begin{equation}\label{eq:old-labels}
 \cR(F_0,D)=\{\ell\in\N:H_i\nmid\ell
                     \text{ for every preceding }H_i\},
\end{equation}
and in particular
\begin{equation}\label{eq:initial-density}
 \dens(\cR(F_0,D))>\frac{99}{100}.
\end{equation}
\end{lemma}

\begin{proof}
Write each $a\in F$ as $A_a/B_a$ with positive integer numerator and
denominator. Replacing $\ell$ by $\ell+A_a$ and $k$ by $k+DB_a$ leaves
$ka-D\ell$ unchanged. Thus the condition for this point is periodic
with period $A_a$, and the product of these finitely many numerators is
a common period. For $F=\varnothing$, every label is admitted and one
may take period $1$.

Equation~\eqref{eq:old-labels} follows from the two alternatives in
Lemma~\ref{lem:grid}. By the Chinese remainder theorem its density is
$\prod_i(1-H_i^{-1})$, with value $1$ for an empty product. The inequality
$\prod_i(1-x_i)\geq1-\sum_i x_i$ for $0\leq x_i\leq1$, together with
\eqref{eq:reciprocal-primes}, proves \eqref{eq:initial-density}.
\end{proof}

The reset in Lemma~\ref{lem:grid} changes only the grid used for future
blocks. Every old source excludes just one class of labels on this grid,
even though all of its integer multiples must be avoided.

\section{The density cost of a block}
\label{sec:cost}

For integers $D\geq1$ and $p\geq2D$, put $q=p-1$, $c=p/q$, and
$r=\lfloor q/D\rfloor$. Define a finite forbidden set of residues by
\begin{equation}\label{eq:bad-residues}
 \mathcal B(p,D)=\{z\bmod p:z\in\Z,\ -r\leq z\leq r\}.
\end{equation}
The interval is closed. This slightly stronger exclusion gives the
half-open admissibility required in \eqref{eq:free-labels}.

\begin{lemma}[Root-residue bound]\label{lem:cost}
Let $\ell\in\N$. If $\ell\bmod p\notin\mathcal B(p,D)$, then
$|cDh-D\ell|>1$ for every $h\in\Z$. Moreover,
\begin{equation}\label{eq:cost}
 \frac{|\mathcal B(p,D)|}{p}\leq\frac3D.
\end{equation}
Consequently the same exclusion separates $D\ell$ from every integer
multiple of every point in any block $\{cDn:n\in P\}$.
\end{lemma}

\begin{proof}
If $|cDh-D\ell|\leq1$, then $D|ph-q\ell|\leq q$.
The integer $z=ph-q\ell$ therefore lies in $[-r,r]$. Since $q=p-1$,
we have $z\equiv\ell\pmod p$, proving the first assertion by
contraposition. Reduction modulo $p$ cannot increase cardinality, so
$|\mathcal B(p,D)|\leq2r+1$. Hence
\[
 \frac{|\mathcal B(p,D)|}{p}
 \leq\frac{2q}{pD}+\frac1p
 \leq\frac2D+\frac1D=\frac3D.
\]
The last estimate already holds for $p\geq D$. For a point $cDn$,
its $k$-fold multiple is $cD(kn)$, so the first assertion applies with
$h=kn$. No coprimality between $D$ and $p$ is assumed.
\end{proof}

\begin{remark}
The number $cD$ is a counting device only. It is not added to the
separated set. Its multiples contain all multiples of the entire block,
which is why the bound is independent of the number of labels in $P$.
\end{remark}

Fix a grid $D$ for the configuration $F_0$ at the start of a round, and
put
\begin{equation}\label{eq:round-length}
 J=\left\lfloor\frac D{12}\right\rfloor.
\end{equation}
After $s\leq J$ insertions, let $\mathcal P$ be the set of the root
numerators $p$ used in this round, so $|\mathcal P|\leq s$ and $p\geq2D$
for $p\in\mathcal P$. Let $\mathcal H$ contain all current private
primes, and possibly a newly chosen one, subject to
$\sum_{H\in\mathcal H}H^{-1}<1/100$. Set
\begin{equation}\label{eq:round-filter}
 E(\mathcal H,\mathcal P,D)=
 \left\{\ell\in\N:
 \begin{array}{l}
 H\nmid\ell\quad(H\in\mathcal H),\\
 \ell\bmod p\notin\mathcal B(p,D)\quad(p\in\mathcal P)
 \end{array}\right\}.
\end{equation}
This filter lies in $\cR(F,D)$ for the entire current point set $F$:
old-round points are protected by Lemma~\ref{lem:grid}, and points
inserted in this round by Lemma~\ref{lem:cost}.

A common period is the product of the finitely many moduli in
$\mathcal H$ and $\mathcal P$, with empty product $1$. For each modulus
$m$ dividing that period, every residue modulo $m$ has the same number
of preimages. Thus lifting a forbidden set to the common period
preserves its density. The finite union bound gives
\begin{equation}\label{eq:round-density}
 \begin{split}
 \dens(E)&\geq1-\sum_{H\in\mathcal H}\frac1H
                 -\sum_{p\in\mathcal P}\frac{|\mathcal B(p,D)|}{p}\\
         &>\frac{99}{100}-\frac{3s}{D}
          \geq\frac{37}{50}>\frac23.
 \end{split}
\end{equation}
No independence or coprimality of these moduli is required. The single
private-prime budget covers both the old-round points and all private
prime exclusions; it is not counted twice.


\section{The next block}
\label{sec:extension}

Retain the round setup of Section~\ref{sec:cost}. The grid $D$ was chosen
for the initial configuration $F_0$ by Lemma~\ref{lem:grid}; the current
point set $F$ consists of $F_0$ and the $s<J$ blocks already inserted on
this same grid. Assume that the current configuration is admissible and
that $D$ is coprime to all its private primes. This holds at the start
of the round by~\eqref{eq:grid-choice}. We show that the next insertion
exists and preserves these conditions; its new private prime will exceed
$D$.

\subsection{Choosing the prime before the labels}
Let $i$ be the index of the next block and let
$\delta=1/100-\sum_{j<i}H_j^{-1}>0$. Choose a prime $H_i$ exceeding
\begin{equation}\label{eq:H-choice}
 2,\quad D,\quad \delta^{-1},\quad
 \text{all preceding }H_j,\quad
 \text{all preceding integers }N_{j,n}.
\end{equation}
Unboundedness of the primes meets these finitely many bounds.
Since $H_i^{-1}<\delta$, the enlarged private-prime set still satisfies
\eqref{eq:reciprocal-primes}. Let $\mathcal H=\{H_1,\ldots,H_i\}$,
let $\mathcal P$ be the set of numerators $p_j$ of the $s$ blocks already
inserted in this round, and define
\begin{equation}\label{eq:E-choice}
 E=E(\mathcal H,\mathcal P,D).
\end{equation}
The entire periodic filter, and a common period, are now fixed. By
\eqref{eq:round-density},
\begin{equation}\label{eq:E-density}
 \dens(E)>\frac{37}{50}>\frac23.
\end{equation}
Also $E\subset\cR(F,D)$ and every label in $E$ avoids every current
private prime, including $H_i$.

Apply Theorem~\ref{thm:selection} to this fixed $E$, with an integer
threshold large enough to obtain a finite primitive $P_i\subset E$ with
\begin{equation}\label{eq:P-choice}
 \begin{gathered}
 W(P_i)>\frac16,\qquad
 \min P_i\geq2D,\\
 D\min P_i>4H_i,\qquad
 D\min P_i>M(F)+1.
 \end{gathered}
\end{equation}
Such a threshold exists because $F,D,H_i$ are already fixed. The set
$P_i$ is nonempty, and we put $U=D\max P_i$.

\subsection{Choosing the perturbation after the labels}
For each $a\in F$, retain the reduced representation $a=A_a/B_a$ with positive
integer denominator $B_a$. Choose a prime $G_i$, distinct from all
$H_1,\ldots,H_i$, sufficiently large that, with
$p=H_iG_i$, $q=p-1$, and $c=p/q$,
\begin{equation}\label{eq:G-choice}
 \begin{gathered}
 p\geq2D,\qquad 1<c<2,\qquad q\geq2U^{13},\\
 \frac Uq<\frac1{B_a}\qquad(a\in F).
 \end{gathered}
\end{equation}
The last family is empty when $F=\varnothing$. All these conditions are
finitely many lower bounds on $G_i$, together with finitely many prime
exclusions, and so can be met simultaneously.

Set $D_i=D$ and define $A_i$ by \eqref{eq:block-data}. For $n\in P_i$,
write $m=Dn$. Then
\begin{equation}\label{eq:eta}
 a_{i,n}=cm=m+\eta_m,\qquad
 \eta_m=\frac mq>0,
\end{equation}
and, as $m\leq U$,
\begin{equation}\label{eq:eta-small}
 \eta_m\leq\frac{m}{2U^{13}}\leq\frac12m^{-12}.
\end{equation}
The choice of $G_i$ changes none of $F,D,H_i,P_i$.

\begin{proposition}[Finite extension]\label{prop:extension}
The enlarged configuration obtained by appending $A_i$ is admissible.
It still has $(D,H_j)=1$ for every private prime currently present, and
the new block satisfies
\begin{equation}\label{eq:block-weight}
 W(A_i)>\frac1{24D}.
\end{equation}
\end{proposition}

\begin{proof}
We first verify every type of ordered pair in the separation condition.

\smallskip
\noindent\emph{An old source and a new target.}
Fix $a=A_a/B_a\in F$, $n\in P_i$, and $k\in\Z$, and set $m=Dn$.
Since $n\in\cR(F,D)$, the number $ka-m$ is outside $(-1,1]$.
It belongs to $B_a^{-1}\Z$, so either
\[
 ka-m\leq-1\qquad\text{or}\qquad ka-m\geq1+\frac1{B_a}.
\]
In the first case $ka-a_{i,n}=ka-m-\eta_m<-1$. In the second,
\eqref{eq:G-choice} gives
\[
 ka-a_{i,n}\geq1+\frac1{B_a}-\eta_m
 >1.
\]
Thus $|ka-a_{i,n}|>1$ for every $k\in\Z$, in particular for every
positive $k$. The denominator $B_a$ is fixed for this source, and the
argument places no upper bound on $k$.

\smallskip
\noindent\emph{A new source and an old target.}
By \eqref{eq:P-choice} and $c>1$, every new point $a_{i,n}$ exceeds
$M(F)+1$. For $a\in F$ and $k\in\N$,
\[
 ka_{i,n}-a\geq a_{i,n}-a>1.
\]

\smallskip
\noindent\emph{Two distinct points in the new block.}
For distinct $n,n'\in P_i$ and $k\in\N$, primitivity implies
$kn\ne n'$. Therefore
\[
 |ka_{i,n}-a_{i,n'}|=cD|kn-n'|\geq cD>D>1.
\]
Old pairs already satisfy (I1). The inequalities above also give the
required ordering of the new block, so (I1) holds.

For (I2), consider first a preceding integer $N_{j,n}$. The old private
primes do not divide it by induction, and the new prime does not divide
it because $H_i>N_{j,n}$. Next consider a new integer
$N_{i,n}=G_iDn$. For every private prime $H_j$ with $j\leq i$:
\begin{itemize}
\item $H_j\nmid D$, by the round's hypothesis for $j<i$ and by $H_i>D$
for $j=i$;
\item $H_j\nmid n$, by $P_i\subset E$;
\item $H_j\nmid G_i$, because $G_i$ is a prime distinct from $H_j$.
\end{itemize}
Since $H_j$ is prime, it does not divide $G_iDn$. This proves
\eqref{eq:private-invariant} for every pair of indices now present.
The size and distinctness requirements on the private primes follow
from \eqref{eq:H-choice}, and the reciprocal budget was preserved when
$H_i$ was selected. No assertion about an unchosen prime is used.

For (I3), equations~\eqref{eq:P-choice} and \eqref{eq:eta-small} give
\[
 \frac{a_{i,n}}{H_i}=\frac{cDn}{H_i}>4,
 \qquad 0<a_{i,n}-Dn\leq\frac12(Dn)^{-12}.
\]
Finally, $n\geq2D$ and $c<2$ imply $cD\leq n$ and hence
$\log(cDn)\leq2\log n$. Consequently,
\begin{equation}\label{eq:weight-comparison}
 \wt(cDn)=\frac1{cDn\log(cDn)}
 \geq\frac{\wt(n)}{4D}.
\end{equation}
The map $n\mapsto cDn$ is injective. Summing
\eqref{eq:weight-comparison} and using $W(P_i)>1/6$ gives
\eqref{eq:block-weight}.
\end{proof}

\section{Rounds and the infinite union}
\label{sec:union}

\begin{proposition}[One round]\label{prop:round}
Every admissible finite configuration can be extended, without changing
any of its points or data, to an admissible finite configuration whose
point set has gained weight at least $1/576$.
\end{proposition}

\begin{proof}
Start with its point set $F_0$ and choose $D$ by Lemma~\ref{lem:grid}.
Keep this $D$ fixed, and prescribe $J=\lfloor D/12\rfloor$ new blocks.
After $s<J$ insertions, inequality~\eqref{eq:round-density} is valid.
Sections~\ref{sec:extension} and \ref{sec:cost} therefore allow the next
insertion: the new configuration is admissible, its $p_i$ satisfies
$p_i\geq2D$, and $D$ remains coprime to every private prime.
Finite induction performs all $J$ insertions.

The blocks are pairwise disjoint by their ordering in (I1). Their total
added weight is therefore, by \eqref{eq:block-weight}, at least
$J/(24D)$. Since $D\geq48$, we have
\[
 J=\left\lfloor\frac D{12}\right\rfloor
 \geq\frac D{12}-1\geq\frac D{24}.
\]
Thus the gain is at least $1/576$. The number of insertions is finite,
even though no size bound on $D$ is imposed. At the end of the round,
(I2) allows Lemma~\ref{lem:grid} to choose a new grid using every integer
$N_{i,n}$ now present. This new grid is used only for future blocks; no
old value $D_i$ is changed.
\end{proof}

\begin{proof}[Proof of Theorem~\ref{thm:main}]
Begin with the empty admissible configuration and apply
Proposition~\ref{prop:round} repeatedly. Let $F_R$ be the point set at the
end of round $R$, and set $F_0=\varnothing$. Each round extends the same
preceding configuration. Thus
\[
 F_0\subset F_1\subset F_2\subset\cdots,
 \qquad W(F_R)\geq\frac R{576}\quad(R\geq1).
\]
Set $A=\bigcup_{R\geq1}F_R$. Every point is rational by
\eqref{eq:block-data}, and every point exceeds $2$, since
$D_i\geq48$ and $n\geq2D_i$ at its insertion.

Given distinct $a,b\in A$, choose a round $R$ with $a,b\in F_R$.
The separation property of $F_R$ yields $|ka-b|\geq1$ for every
$k\in\N$. This proves \eqref{eq:original} on the union, since the
round $R$ containing the pair is chosen independently of $k$.

The set $A$ is a countable union of finite sets. Its separation at $k=1$
implies that every bounded interval contains only finitely many points,
so $A$ has no finite accumulation point. Equation~\eqref{eq:near-integer}
holds for every point by (I3), and remains true because no point is ever
changed.

Finally, monotonicity of nonnegative sums gives
\[
 W(A)\geq W(F_R)\geq\frac R{576}\qquad\text{for every }R\geq1.
\]
Therefore $W(A)=\infty$. Since each point of $A$ has finite weight,
$A$ must be infinite. This proves all assertions of the theorem.
\end{proof}

\section{Order of choices and formal verification}
\label{sec:status}

\subsection{Quantifier order}
The dependency order for a new block is
\[
 \text{preceding configuration}
 \ \longrightarrow\ D\ \longrightarrow\ H_i\ \longrightarrow\ E
 \ \longrightarrow\ P_i\ \longrightarrow\ G_i\ \longrightarrow\ A_i.
\]
The grid $D$ is fixed throughout a round. Every use of
Theorem~\ref{thm:selection} fixes the entire periodic set $E$, including
all prime exclusions, before selecting $P_i$. Only then is $G_i$ chosen.
The analytic limit therefore has a fixed modulus. Compatibility with
later primes is preserved by choosing each new $H_i$ above all stored
$N_{j,n}$, while the new labels and $G_i$ avoid the current private primes.

\subsection{Choice of successive extensions}
The formal construction selects a witness of the proved finite-round
extension theorem by classical choice and iterates this selection by
recursion on the natural numbers. It establishes existence without a
quantitative bound on the sizes of the grids, periods or blocks, or an
extracted executable construction.

\subsection{Verification scope and reproducibility}
\label{sec:verification}
The complete existence theorem is formalized in Lean~4.24.0. All names
below are in the namespace \path{Erdos143}. The two closed rational
theorems are
\begin{center}\small
\begin{tabular}{@{}l@{\quad}l@{}}
\path{mainConclusion_proved} & \path{MainConclusion}\\
\path{strongConclusion_proved} & \path{StrongConclusion}
\end{tabular}
\end{center}
The first gives a countable rational set in $[2,\infty)$, separation for
every ordered distinct pair and every positive natural multiplier, and
infinite nonnegative total weight. The second adds local finiteness and
\eqref{eq:near-integer}. Their definitions are in
\path{Erdos143/Definitions.lean} and their proofs in
\path{Erdos143/Main.lean}. The project proves the periodic selection lemma
and the finite-round extension theorem, then constructs the nested
sequence using classical choice as described above.

The real-domain statement is proved in \path{RealConclusion.lean} as
\begin{center}\small
\path{Erdos143.realConclusion_proved}.
\end{center}
It explicitly includes countability,
rationality, the lower bound $2$, every ordered multiplier inequality,
and both divergence of the nonnegative sum and nonsummability of the
real-valued weight. The transfer uses the injective embedding of
$\Q$ into $\R$ and an equivalence of the indexing sets. In particular,
divergence is encoded by $\top$ in the extended nonnegative reals
(Lean's \path{ENNReal}), together with an explicit negation of
\path{Summable}; it is not inferred from the conventional value of a
nonsummable real \path{tsum}. The lower bound and divergence imply that
the set is infinite.

The periodic selection theorem is checked by the declaration
\begin{center}\small
\path{periodicPrimitiveSelection_proved}.
\end{center}
The finite construction follows the same order of choices and uses the
same strict prime budget, closed residue filter and product reset grid
as the proof above. The declarations
\begin{center}\small
\begin{tabular}{@{}l@{\quad}l@{}}
\path{insert_block_exists} & finite insertion\\
\path{block_extension_valid} & simultaneous invariants\\
\path{round_extension_exists} & gain of at least $1/576$\\
\path{stateSequence_mass_lower} & prefix weight bound
\end{tabular}
\end{center}
connect these steps to the nested union in \path{Main.lean}.
The formal states store finite sets of triples $(H,N,q)$ and private
primes; the prose groups the triples into blocks. A detailed
statement-by-statement correspondence is supplied with the code in the
\path{docs/} directory. No asymptotic assertion for every individual
$\Om$ level is used. The supplementary file
\path{ManuscriptAudit.lean} also checks a single existential statement
combining these conclusions, including infinitude and the same nested
finite prefixes.

\Needspace{5\baselineskip}
The environment pins mathlib to revision
\begin{center}\small
\texttt{f897ebcf72cd16f89ab4577d0c826cd14afaafc7},
\end{center}
with all transitive revisions fixed in \path{lake-manifest.json}.

\Needspace{7\baselineskip}
From the supplied supplementary project, the commands
\begin{quote}\small\ttfamily
bash scripts/bootstrap.sh\\
bash scripts/verify.sh
\end{quote}
retrieve the pinned dependencies, build the proof with warnings treated
as errors, audit all proved declarations and compile the expanded
statement checks. Network access and the tools listed in the project
\path{README.md} are required. On platforms other than Linux x86\_64,
Lean~4.24.0 must first be installed using the platform's standard elan
installation.

The supplementary material contains the two original local verification
records and a further build and statement audit dated October 23, 2026.
The original runs include a separate-copy clean rebuild of the project
using pinned dependency caches. These checks do not constitute a
from-source rebuild of Lean and mathlib.
All 245 proved declarations passed the axiom audit. The three closed
conclusion theorems use only \path{propext}, \path{Classical.choice} and
\path{Quot.sound}, the standard foundations of the Lean/mathlib
formalization; no added analytic or arithmetic axiom is used.
The original frozen source archive, its SHA-256 manifest and the local
logs are included in the supplementary project's \path{verification/}
directory. Typesetting this article does not execute the formal proof.
Kernel checking establishes the stated formal results from their recorded
foundations. Independent expert assessment of the formulation, its
correspondence to Erd\H{o}s's question and the exposition remains part
of the journal's review process.

\begin{thebibliography}{9}

\bibitem{Bloom143}
T. F. Bloom, \emph{Erd\H{o}s problem \#143}, Erd\H{o}s Problems,
\url{https://www.erdosproblems.com/143}, problem formulation consulted October 2, 2026.

\bibitem{Erdos1935}
P. Erd\H{o}s, \emph{Note on sequences of integers no one of which is
  divisible by any other}, J. London Math. Soc. \textbf{10} (1935), 126--128.
\href{https://www.renyi.hu/~p_erdos/1935-04.pdf}{Author archive copy}.

\bibitem{Erdos1948}
P. Erd\H{o}s, \emph{On the density of some sequences of integers},
Bull. Amer. Math. Soc. \textbf{54} (1948), 685--692, especially problem II,
p.~692.
\href{https://doi.org/10.1090/S0002-9904-1948-09057-7}{doi:10.1090/S0002-9904-1948-09057-7}.
\href{https://www.renyi.hu/~p_erdos/1948-06.pdf}{Author archive copy}.

\bibitem{KLL2025}
D. Koukoulopoulos, Y. Lamzouri, and J. D. Lichtman,
\emph{Erd\H{o}s's integer dilation approximation problem and GCD graphs},
arXiv:2502.09539v1 (2025), Theorem~1.
\url{https://arxiv.org/abs/2502.09539v1}.

\bibitem{Tao2014}
T. Tao, \emph{254A, Notes 1: Elementary multiplicative number theory},
What's New, November 23, 2014.
\href{https://terrytao.wordpress.com/2014/11/23/254a-notes-1-elementary-multiplicative-number-theory/}{Online lecture notes}.

\end{thebibliography}
\end{document}
